\documentclass[10pt,a4paper]{amsart}
\usepackage[margin=19mm]{geometry}
\usepackage{amsmath,amssymb,mathtools}
\usepackage[scr=rsfs,cal=euler]{mathalfa}
\usepackage{xfrac}
\usepackage[unicode,hidelinks,bookmarks=false]{hyperref}
\newcommand{\ie}{i.e.\ }
\newcommand{\cf}{cf.\ }
\newcommand{\resp}{respectively\ }
\newcommand{\Subsection}{\subsection}
\providecommand{\nobreakdash}{\nobreak\hbox{-}}
\setlength{\emergencystretch}{2em}
\multlinegap0pt
\usepackage{bm}
\usepackage{bbm}
\usepackage{dsfont}
\usepackage{xspace}
\usepackage{cancel}
\usepackage{mathtools}
\usepackage[shortlabels]{enumitem}
\usepackage{amsthm}
\makeatletter
\@ifundefined{theorem}{\newtheorem{theorem}{Theorem}}{}
\@ifundefined{lemma}{\newtheorem{lemma}[theorem]{Lemma}}{}
\makeatother
\newcommand{\F}{\mathcal{F}}
\newcommand{\fast}{{\sc Fast}}
\newcommand{\gcg}{\gamma_{\rm cg}}
\title[Complexity of the game connected domination problem]{Complexity of the game connected domination problem}
\author{Vesna Ir\v{s}i\v{c} Chenoweth}
\date{Source: arXiv:2404.08776v2 (CC BY 4.0)}
\begin{document}
\maketitle

\noindent\emph{Source and licence.} Complexity of the game connected domination problem. Version arXiv:2404.08776v2 (CC BY 4.0), distributed under the Creative Commons Attribution 4.0 International licence, as recorded in the arXiv licence metadata for this exact version and shown on the abstract page https://arxiv.org/abs/2404.08776v2. Full rights notice: licenses/arxiv-2404.08776-CC-BY-4.0.txt.\par\smallskip

\noindent\textit{Editorial scope.} Counted unit: the Lemma whose source label is \texttt{lem:p1-wins} in the article (source0.tex lines 672--678, source label \texttt{lem:p1-wins}), delivered together with the article's own complete original proof of it (source0.tex lines 680--746, one \texttt{proof} environment of 67 lines). No mathematics is written, repaired or re-derived: statement and proof are the article's own text line for line. Required context retained uncounted: none. Every definition, display and label of those lines is retained unchanged. Omitted from this package and not redistributed: 1--671, 679--679, 747--1338. Nothing inside a retained excerpt is omitted or altered: every retained body is delivered exactly as the article's lines are written, so no omission receipt of any kind accompanies this package. Citations: the retained excerpts carry no citation command at all, so no bibliography entry of the article is transcribed and no reference is redistributed. Reference adaptation: none was needed and none was made; every reference inside the delivered unit resolves inside it, so no unresolved reference or citation is produced. Printed slips kept literal: no printed word, symbol, hypothesis or number of the counted statement or of its proof was altered. Layout and class: the article's own class is \texttt{article} with packages including babel, amsmath, amssymb, amsfonts, amsthm, tikz, color, pgf, graphicx, url, enumerate, xcolor, hyperref, geometry; the wrapper is the lane's pinned \texttt{amsart} editorial wrapper at 10pt on A4, which restates the theorem-like environments the retained text needs and adds the article's own macro definitions that the retained text needs (\texttt{\textbackslash F}, \texttt{\textbackslash fast}, \texttt{\textbackslash gcg}), with extra wrapper packages (bm, bbm, dsfont, xspace, cancel, mathtools, enumitem). The delivered numbering is the unit's own. Assets: no retained span contains a figure environment, a graphics-inclusion command, a PDF-page inclusion, a table or a TikZ picture; no image file is copied into this package, the package declares no asset dependency and no image file is redistributed with it. Licence: version arXiv:2404.08776v2 of this article is distributed under the Creative Commons Attribution 4.0 International licence, as recorded in the arXiv licence metadata for this exact version and shown on the abstract page https://arxiv.org/abs/2404.08776v2. A full rights notice is delivered at licenses/arxiv-2404.08776-CC-BY-4.0.txt. Characters: every retained excerpt is pure ASCII, so the delivered engine, XeLaTeX with its default Latin Modern text font, reports no missing character.
\begin{lemma}
	\label{lem:p1-wins}
	If Player 1 has a winning strategy for the POS-CNF game played on $\F$, then $$\gcg(G_\F) \leq \begin{cases}
		3k + 4n + 8; & k \text{ even},\\
		3k + 4n + 11; & k \text{ odd}.
	\end{cases}$$
\end{lemma}

\begin{proof}
	We describe Dominator's strategy that ensures that the game ends in at most $3k + 4n + 8$ moves if $k$ is even and in at most $3k + 4n + 11$ moves if $k$ is odd. Most of the proof is the same for both parities of $k$.
	
	Dominator starts the game by playing $d_1 = u_{2n+7}$. Thus he forces the strategy \fast\ to be played on $H_{2n+7}$, so the moves $s_1 = u_{2n+6}, d_2 = u_{2n+5}, \ldots, d_{2n+7} = u_1, s_{2n+8} = u_0$ are forced. As $2n+8$ is even, Staller plays $u_0$. Observe that no vertex $d_j^m$ can be a legal move in the rest of the game.
	
	Dominator's next move is to play $d_{2n+9} = a_i$ if setting $X_i$ to TRUE is the optimal first move of Player 1 in the POS-CNF game on $\F$. For the rest of the game, Dominator will imagine a game is being played on $\F$, specifying certain moves of Staller and himself as moves in the POS-CNF game. More precisely, in the game on $\F$, Player 1 will be playing optimally, and their moves will determine some of the corresponding moves of Dominator on $G_\F$. As Player 1 has a winning strategy on $\F$, they can win no matter how Player 2 plays. Player 2's moves on $\F$ will be determined by some of Staller's moves on $G_\F$. Which moves of Staller mean that Player 2 makes a move on $\F$ and how does Player 1's reply on $\F$ translate to Dominator's reply on $G_\F$ is explained in (6) below.
	
	We will prove that Dominator can ensure on average at most three moves on each $B_i$, at most two moves on each $C_j$ and at most three moves on $A$. By saying on average we mean that the moves can be counted as at most three on each $B_i$ and at most two on each $C_j$, as Dominator's strategy ensures that if four moves are played on some $B_i$, there is a $C_j$ with only one move played on it. 
	
	Dominator follows the rules listed below (each rule is roughly summarized first, followed by the exact Dominator's strategy) with the following addition: when we say that Dominator plays \emph{any legal move}, we mean that he plays any $c_j$, $a_i$ or $e_i$ if possible (in this order of preference), and otherwise he plays any vertex. This is to ensure at most two moves on each $C_j$ and at most three moves on each $B_i$ even when Dominator does not have a more precise strategy. As a result, it might sometimes happen that Dominator's strategy below wants him to play a vertex that is not a legal move anymore. In this case, he plays any legal move (with the rules given above), but he still performs potential other actions associated with the move that is not legal anymore as if he played it now. It is thus necessary for Dominator to keep track of the vertices he played due to an exact strategy and the vertices he played as any legal moves. 
%	\begin{enumerate}[(1)]
%		\item If Staller plays $x_1$, Dominator replies by playing $x_2$.
%		\item If Staller plays $x_3$ or $y_2$, then Dominator plays any legal move.
%		\item If Staller plays on $C_j$, Dominator plays $c_j$ if possible.
%		\item If Staller plays on $B_i$ and not all vertices of $B_i$ are dominated after her move, Dominator replies on $B_i$ as well.
%		\item If Staller plays $a_i$ and all vertices of $B_i$ are dominated already before her move, Dominator replies by playing $c_j$ in the appropriate $C_j$.
%		\item If Staller plays on $B_i$ that was not dominated before her move, all vertices of $B_i$ are dominated after her move, and not all gadgets have been played on yet, Dominator selects an appropriate gadget $B_i$ and plays the vertex $a_i$ there.
%		\item Otherwise, Dominator plays some $c_j$ or some $e_i$ or any legal move (in this order of preference).
%	\end{enumerate}
%	A more detailed description of Dominator's strategy in each of these cases is given below:	
	\begin{enumerate}[(1)]
		\item If Staller plays $p_1$, Dominator replies by playing $p_2$.\\
		This ensures that if Staller is the first to play on $A$, then exactly three moves are made on $A$ during the game.
		
		\item If Staller plays $p_3$ or $q_2$, then Dominator plays any legal move.\\
		Note that this means that on $G - H_{2n+7} -A$, Dominator makes two consecutive moves which we have to consider when analyzing some of the cases below.
		
		\item If Staller plays on $C_j$, Dominator plays $c_j$.\\
		If Dominator played first on $C_j$ (before this move of Staller), then his strategy ensures he played $c_j$, thus no more moves are legal on $C_j$. So we know that the current move of Staller is in fact the first move on $C_j$. If there are still undominated vertices on $C_j$ after Staller's move, $c_j$ is a legal reply for Dominator that ensures that at most two moves are made on $C_j$ during the game. Otherwise, so if Staller's move on $C_j$ leaves no undominated vertices in $C_j$, then Staller played $c_j$, Dominator plays any legal move, and just one move is made on $C_j$ during the game.
		
		\item If Staller plays on $B_i$ and not all vertices of $B_i$ are dominated after her move, Dominator replies on $B_i$ as well.\\
		More precisely, if Staller is the first to play on some $B_i$, then Dominator replies by playing $e_i$. If Staller is not the first to play on $B_i$, but $B_i$ is still not entirely dominated, then this is only possible if Dominator played the first move on this $B_i$, so by (6), he played $a_i$ before. If Staller now plays $e_i$, he replies on $b_i$, and if she plays $b_i$, he selects $e_i$ (no other move on $B_i$ is a legal move for Staller at this stage of the game). Note that if Dominator made two consecutive moves on $B_i$, they were $a_i$ and $e_i$, thus after Staller's move on $B_i$, there are no more undominated vertices left and we are not in this case. Similarly if Dominator made three consecutive moves on $B_i$.
		
		Notice that with this strategy, Dominator ensures that no more than three moves are played on each $B_i$ unless Staller plays a move from (5).
		
		\item If Staller plays $a_i$ and all vertices of $B_i$ are dominated already before her move, Dominator replies by playing $c_j$ in the appropriate gadget $C_j$.\\
		Note that if such a move of Staller was legal, it must have newly dominated vertices $c_j, c_j^1, \ldots, c_j^n$ for some $j \in [n]$. Thus Dominator's reply on $c_j$ is a legal move. Additionally, notice that while four moves were played on $B_i$, only one move was played on $C_j$ (and no more moves on $B_i$ or $C_j$ are possible during the game). For counting purposes, we consider this as three moves on $B_i$ and two on $C_j$.
		
		\item If Staller plays on $B_i$ that was not dominated before her move and all vertices of $B_i$ are dominated after her move, Dominator selects an appropriate gadget $B_{j}$ to play on.\\
		If $X_i$ already has an assigned value in the imagined POS-CNF game on $\F$, then Dominator plays any legal move. Otherwise, Dominator imagines that Player 2 set $X_i$ to FALSE in the POS-CNF game on $\F$ now. If this ends the POS-CNF game on $\F$, Dominator plays any legal move. Otherwise, Dominator considers Player 1's optimal strategy in the POS-CNF game, and plays on $B_{j}$ such that Player 1's strategy is to set $X_{j}$ to TRUE. Dominator's strategy is to play $a_j$ if possible, otherwise he plays $b_j$ if this is the last move on $B_j$, or he plays any legal move (in this order of preference). Once Dominator sets $X_j$ to TRUE as Player 1's move on $\F$ and wants to play on $B_j$, the following options are possible for $B_j$: no move has been played on it so far (Dominator now plays $a_j$), Dominator already played one or more moves on it (Dominator now plays $b_j$ if this is the last move on $B_j$, or he plays any legal move, but we know he played $a_j$ before), Staller and Dominator both played on $B_j$ before (Dominator now plays $a_j$ if possible, or $b_j$ if this is the last move on $B_j$, or he plays any legal move, but again we know that $a_j$ has been played before). In all cases we notice that Dominator is ensuring that at most three moves have been or will be played on $B_j$, $a_j$ is one of them, and that if Staller was the first to play on $B_j$, all vertices of $B_j$ are dominated after this move.
		
%		If Staller plays on $B_i$ that was not dominated before her move, all vertices of $B_i$ are dominated after her move, and at least two variables have not yet been assigned values in the POS-CNF game, Dominator selects an appropriate gadget $B_{j}$ to play on.\\
%		In this case, Staller played the third move on some gadget $B_i$, and Dominator imagines that Player 2 set $X_i$ to FALSE in the POS-CNF game on $\F$. Note that by (4) this only happens if Staller also played the first vertex of $B_i$. Dominator now considers Player 1's optimal strategy in the POS-CNF game, and plays on $B_{j}$ such that Player 1's strategy is to set $X_{j}$ to TRUE. Dominator's strategy is to play $a_j$ if possible, otherwise he plays $b_j$ if this is the last move on $B_j$, or he plays any legal move. Notice that as $X_j$ has not been set to FALSE previously, if three moves have been played on $B_j$ before, the last move was not made by Staller and Dominator played $a_j$ before. Once Dominator sets $X_j$ to TRUE and wants to play on $B_j$, the following options are possible for $B_j$: no move has been played on it so far (Dominator now plays $a_j$), Dominator already played one or more moves on it (Dominator now plays $b_j$ if this is the last move on $B_j$, or he plays any legal move), Staller and Dominator both played on $B_j$ before (Dominator now plays $b_j$ if this is the last move on $B_j$, or he plays any legal move). In all cases we notice that Dominator is ensuring that at most three moves have been or will be played on $B_j$, and that if Staller was the first to play on $B_j$, all vertices of $B_j$ are dominated after this move.
		
%		If no moves have been played on $B_j$ so far, Dominator plays $a_j$. If one move has been played on $B_j$ so far, it was the move of Domiantor by (4), so we know he played $a_j$ before. Now Dominator plays any legal move. If two moves have been played on $B_j$ so far, we know by (4) that it was first the move of Staller and then Dominator's reply. Thus either $a_j, e_j$ have been played so far, in this case, Dominator plays $b_j$, or $b_j, e_j$ have been played so far, in this case Dominator plays $a_j$ now. If three moves have been played on $B_j$ already, it cannot be that Staller played the first and last move (by the first part of (6)), so Dominator played the first and last on $B_j$, so $a_j$ has been played already, and now Dominator plays any legal move.
		
		Notice that with this strategy, Dominator ensures that for all indices $i \in I$ that Player 1 sets to TRUE in his winning strategy in the POS-CNF game on $\F$, the vertices $a_i$, $i \in I$, are played. Thus by the end of the POS-CNF game, all vertices $c_1, \ldots, c_n, \ldots, c_1^n, \ldots, c_n^n$ will be dominated in $G_\F$ (since Player 1 wins on $\F$) and so all vertices from $C_1, \ldots, C_n$ can be dominated by using at most three moves on each $B_i$ and at most two on each $C_j$ (on average).
		
%		If Staller plays on $B_i$ that was not dominated before her move, all vertices of $B_i$ are dominated after her move, and not all gadgets have been played on yet, Dominator selects an appropriate gadget $B_{i'}$ and plays the vertex $a_{i'}$ there.\\
%		In this case, Staller played the third move on some gadget $B_i$, and Dominator imagines that Player 2 set $X_i$ to FALSE in the POS-CNF game. Note that by (4) this only happens if Staller also played the first vertex of $B_i$. Dominator now plays on a previously unplayed gadget. He considers Player 1's optimal strategy in the POS-CNF game, and plays $a_{i'}$ such that Player 1's strategy is to set $X_{i'}$ to TRUE.
%		
%		Notice that with this strategy, Dominator ensures that for all indices $i \in I$ that Player 1 sets to TRUE in his winning strategy in the POS-CNF game on $\F$, the vertices $a_i$, $i \in I$, are played. This results in the fact that by the end of the POS-CNF game, all vertices $c_1, \ldots, c_n, \ldots, c_1^k, \ldots, c_n^k$ will be dominated in $G_\F$ (since Player 1 wins on $\F$).
		
%		\item Otherwise, Dominator plays any legal move.\\
%		As we are in this case, the imaginary POS-CNF game has ended already with the win of Player 1, so all vertices $c_1, \ldots, c_n, \ldots, c_1^k, \ldots, c_n^k$ are dominated. If there are any vertices left undominated on some $B_i$ (by above, Dominator played $a_i$ in such gadgets), then Dominator plays $e_i$ (if possible), or $b_i$ if this will be the last move on $B_i$. Again, this strategy ensures that at most three moves are played on each gadget $B_i$. If the only undominated vertices left are in gadgets $C_j$, then Dominator plays $c_j$ in any of them. By (3), this is the first and only move on $C_j$.
%		
%%		If $X_i$ does not yet have a value set in the POS-CNF game, Dominator sets $X_i$ to FALSE as the move of Player 2 in the POS-CNF game.
%%		Now the imaginary POS-CNF game has ended already with the win of Player 1, so all vertices $c_1, \ldots, c_n, \ldots, c_1^k, \ldots, c_n^k$ are dominated. If there are any vertices left undominated on some $B_i$ (by above, Dominator played $a_i$ in such gadgets), then Dominator plays $e_i$ (if possible), or $b_i$ if this will be the last move on $B_i$. Again, this strategy ensures that at most three moves are played on each gadget $B_i$. If the only undominated vertices left are in gadgets $C_j$, then Dominator plays $c_j$ in any of them. By (3), this is the first and only move on $C_j$.
%		
%		Last, we describe Dominator's strategy if $k$ is odd and only vertices in $A$ are undominated. If $p_1$ has been played before, Dominator plays $p_2$ or $p_3$, ensuring at most three moves are made on $A$. If $p_1$ has not been played yet and it is Dominator's turn now, Dominator must play $p_1$, thus possibly four moves are made on $A$. (His consequent move is $p_2$ or $p_3$.) If at most $2n+8+3k+2n-1 = 4n+3k+7$ moves were made during the game before this move, the game ends after at most $4n+3k+11$ moves. Otherwise, so if exactly $2n+8+3k+2n = 4n+3k+8$ moves were made before (by arguments above Dominator ensures at most three moves on each $B_i$ and at most two on each $C_j$), as $k$ is odd, $4n+3k+8$ is also odd, so it cannot be Dominator's turn now.
	\end{enumerate}
	With the described strategy, Dominator ensures exactly $2n+8$ moves on $H_{2n+7}$, and on average at most three moves on each $B_i$ and at most two moves on each $C_j$. Observe also that the imagined POS-CNF game either ends with the win of Player 1 (meaning that all vertices $c_1, \ldots, c_n, \ldots, c_1^n, \ldots, c_n^n$ are dominated and are thus legal moves can be played on any $C_1, \ldots, C_n$ that is not dominated yet) or the POS-CNF game does not end as some variables never get assigned values. By (6) we know that Dominator played the third and last move on all the remaining (unassigned) variables. If the first player playing on such $B_i$ did no play $a_i$, then we know Staller made the move on $b_i$ and Dominator replied by playing $e_i$. Thus by Dominator's strategy, he can play $a_i$ as the last move on $B_i$. If Dominator now finished playing the POS-CNF game with the remaining variables, selecting moves of Player 2 arbitrarily, he would still win, and by the argument above, it still holds that for all indices $i \in I$ that Player 1 sets to TRUE in the POS-CNF game on $\F$, the vertices $a_i$, $i \in I$, are played. Hence, in this case, all vertices $c_1, \ldots, c_n, \ldots, c_1^n, \ldots, c_n^n$ are dominated as well. Thus Dominator also ensures that by keeping at most three moves played on each $B_i$ and at most two on each $C_j$, all vertices from $G_\F - H_{2n+7} - A$ can be dominated (i.e.\ it cannot happen that no move on some $C_j$ would be legal). 
	
	If $k$ is even, there are no moves on $A$. If $k$ is odd and Staller made the first move on $A$, then by (1), three moves are played on $A$. Thus the only remaining case to consider is if $k$ is odd and Dominator is forced to play the first move on $A$. 	This happens if only vertices on $A$ are undominated, $p_1$ has not been played before, and it is Dominator's turn now, Dominator must play $p_1$ (it is the only legal move in the game), thus possibly four moves are made on $A$. (His consequent move is $p_2$ or $p_3$.) If at most $2n+8+3k+2n-1 = 4n+3k+7$ moves were made during the game before this move, the game ends after at most $4n+3k+11$ moves. Otherwise, so if exactly $2n+8+3k+2n = 4n+3k+8$ moves were made before (by arguments above Dominator ensures at most three moves on each $B_i$ and at most two on each $C_j$), as $k$ is odd, $4n+3k+8$ is also odd, so it cannot be Dominator's turn now. Thus if four moves are played on $A$, then at most $4n+3k+7$ moves were played on the rest of the graph.

	It follows from the above that the game ends in at most $2n+8+3k+2n = 4n+3k+8$ moves if $k$ is even, and in at most $2n+8+3k+2n+3=4n+3k+11$ moves if $k$ is odd.	
\end{proof}

\end{document}
